% Rewritten end to end.
%
% Previously every slide in this file was a screenshot -- the method-flow
% diagram and the DDPG/DQN pseudocode from the papers -- with red rectangles
% positioned by hard-coded fractions of the image, and all of it in the paper's
% mu(s|theta^mu) / Q(s,a|theta^Q) notation. Both the diagram and the algorithm
% are now drawn natively in course notation (theta = actor, phi = critic,
% superscript minus = target, D = replay buffer), and the red boxes are attached
% to the content rather than to pixel coordinates.

% ---------------------------------------------------------------------------
% The DDPG data-flow diagram, built in three stages.
% #1 = stage (1 = acting only, 2 = + replay buffer, 3 = + learning loop)
% ---------------------------------------------------------------------------
\newcommand{\ddpgflow}[1]{%
\begin{tikzpicture}[
    font=\scriptsize,
    ell/.style={ellipse, draw=blue!45!black, fill=blue!22,
                inner sep=3pt, align=center},
    buf/.style={rectangle, fill=orange!80!red, text=white, minimum width=2.6cm,
                minimum height=0.66cm, align=center, font=\scriptsize\bfseries},
    net/.style={rectangle, fill=yellow!75, draw=yellow!55!black,
                minimum width=2.3cm, minimum height=1.35cm, align=center},
    tgt/.style={rectangle, fill=green!35, draw=green!45!black,
                minimum width=1.3cm, minimum height=0.75cm, align=center},
    use/.style={rectangle, fill=blue!32, draw=blue!50!black,
                minimum width=1.7cm, minimum height=1.05cm, align=center},
    ar/.style={-{Latex[length=2mm]}, blue!55!black, line width=0.75pt},
  ]
  % Pin the bounding box so the three stages do not jump around.
  \useasboundingbox (-2.0,-1.15) rectangle (13.5,5.0);

  % --- Stage 1: acting in the environment --------------------------------
  \node[net]  (actor)    at (1.5, 1.15) {$\mu_\theta(s)$\\[0.25em] DNN};
  \node[tgt]  (actortgt) at (0.55,-0.45) {$\mu_{\theta^{-}}$\\ slow};
  \node[ell]  (env)      at (-1.0, 3.1) {$s_t$};
  \node[ell]  (trans)    at (4.0, 4.35) {$(s_t,\, a_t,\, r_t,\, s_{t+1})$};

  \draw[ar] (env) -- (actor.130);
  \draw[ar] (actor.55) -- node[left, pos=0.75, font=\scriptsize] {action} (trans.230);
  \node[font=\scriptsize, anchor=west] at (2.55, 1.95) {$+\,\mathcal{N}_t$};

  \ifnum#1=1
    \draw[red, dashed, line width=0.9pt] (-1.95,-0.95) rectangle (5.9, 4.9);
    \node[red, font=\scriptsize\bfseries, anchor=north west] at (6.2, 4.6)
         {act with the actor, $+$ noise};
  \fi

  % --- Stage 2: the replay buffer ----------------------------------------
  \ifnum#1>1
    \node[buf] (buffer) at (9.2, 4.35) {Replay buffer $\mathcal{D}$};
    \draw[ar] (trans) -- (buffer);
    \node[font=\scriptsize, anchor=south] at (6.6, 3.32)
         {Sample random mini-batches};
    \draw[blue!55!black, line width=0.75pt] (buffer.south) -- (9.2, 3.25);
    \draw[blue!55!black, line width=0.75pt] (4.5, 3.25) -- (11.4, 3.25);
  \fi
  \ifnum#1=2
    \draw[red, dashed, line width=0.9pt] (1.9, 3.1) rectangle (11.0, 4.9);
    \node[red, font=\scriptsize\bfseries, anchor=north west] at (2.0, 2.95)
         {store every transition --- and reuse it};
  \fi

  % --- Stage 3: the learning loop ----------------------------------------
  \ifnum#1>2
    \node[align=center] (pgform) at (4.5, 2.62)
         {$\nabla_\theta \mathcal{J} \approx \nabla_\theta \mu_\theta\, \nabla_a Q_\phi$};
    \node[align=center] (yform)  at (11.4, 2.62)
         {$y_i = r_i + \gamma\, Q_{\phi^{-}}\!\big(s_i',\, \mu_{\theta^{-}}(s_i')\big)$};
    \draw[ar] (4.5, 3.25)  -- (pgform.north);
    \draw[ar] (11.4, 3.25) -- (yform.north);

    \node[use] (pgbox)  at (4.5, 1.15)  {use as\\ policy\\ gradient};
    \node[use] (tgtbox) at (11.4, 1.15) {use as\\ target\\ value};
    \node[net] (critic) at (8.5, 1.15)  {$Q_\phi(s,a)$\\[0.25em] DNN};
    \node[tgt] (critictgt) at (7.55,-0.45) {$Q_{\phi^{-}}$\\ slow};

    \draw[ar] (pgform.south)  -- (pgbox.north);
    \draw[ar] (yform.south)   -- (tgtbox.north);
    \draw[ar] (pgbox.west)    -- (actor.east);
    \draw[ar] (tgtbox.west)   -- (critic.east);

    % Same region the original slide boxed: everything except the raw state.
    \draw[red, dashed, line width=0.9pt] (0.1,-0.95) rectangle (13.0, 4.9);
  \fi
\end{tikzpicture}}

% ---------------------------------------------------------------------------
\begin{frame}{Method: The DDPG Loop (1/3) --- Acting}
\vspace{0.2em}
\centering
\resizebox{0.95\textwidth}{!}{\ddpgflow{1}}

\vspace{0.5em}
\begin{flushleft}\footnotesize
The actor is a plain function $s \mapsto \mu_\theta(s)$. Being deterministic it would repeat itself forever, so exploration noise $\mathcal{N}_t$ is added before acting. Each step yields one transition.
\end{flushleft}
\end{frame}

\begin{frame}{Method: The DDPG Loop (2/3) --- Storing}
\vspace{0.2em}
\centering
\resizebox{0.95\textwidth}{!}{\ddpgflow{2}}

\vspace{0.5em}
\begin{flushleft}\footnotesize
Transitions go into the replay buffer $\mathcal{D}$ rather than being used once and discarded --- the DQN trick from Day 2, and the reason DDPG needs far less environment interaction than PPO on the same task.
\end{flushleft}
\end{frame}

\begin{frame}{Method: The DDPG Loop (3/3) --- Learning}
\vspace{0.2em}
\centering
\resizebox{0.95\textwidth}{!}{\ddpgflow{3}}

\vspace{0.3em}
\begin{flushleft}\footnotesize
One mini-batch drives \emph{both} updates: on the right the ``slow'' target networks build the TD target $y_i$ for the critic; on the left the critic's action-gradient becomes the actor's policy gradient. Nothing inside the dashed box asks which policy collected the data --- that is what off-policy means.
\end{flushleft}
\vspace{-0.4em}
\begin{flushright}\tiny Diagram adapted from \rref{9}\end{flushright}
\end{frame}

% ---------------------------------------------------------------------------
% The algorithm, boxing the same regions the paper screenshot used to
% highlight. One frame per highlight rather than beamer overlays: algorithm2e's
% environment is a float, and a float inside a multi-overlay frame breaks
% hyperref's per-slide page anchors (\Hy@PageAnchorSlide).
%
% #1 selects which region is boxed:
%   0 = none, 1 = replay buffer + soft target update, 2 = exploration noise,
%   3 = value (critic) network update, 4 = policy (actor) network update.
% ---------------------------------------------------------------------------
\newcommand{\hlif}[3]{\ifnum#1=#2 \hlbox{red}{#3}\else\hlbox{none}{#3}\fi}

\newcommand{\ddpgalgo}[1]{%
\scriptsize
\begin{algorithm}[H]
\DontPrintSemicolon
\caption{DDPG --- Deep Deterministic Policy Gradient}
Initialise critic $Q_\phi(s,a)$ and actor $\mu_\theta(s)$ with random weights $\phi,\ \theta$\;
Initialise target weights $\phi^{-} \leftarrow \phi$ and $\theta^{-} \leftarrow \theta$\;
\hlif{#1}{1}{Initialise replay buffer $\mathcal{D}$}\;
\For{$\mathrm{episode} = 1$ \KwTo $M$}{
    Observe initial state $s$\;
    \For{$t = 1$ \KwTo $T$}{
        \hlif{#1}{2}{Select action $a = \mu_\theta(s) + \mathcal{N}$ \quad \emph{(actor output $+$ exploration noise)}}\;
        Execute $a$; observe reward $r$, next state $s'$, terminal flag $d$\;
        \hlif{#1}{1}{Store transition $(s,a,r,s',d)$ in $\mathcal{D}$; \quad $s \leftarrow s'$}\;
        \hlif{#1}{1}{Sample a minibatch of $N$ transitions $(s_i,a_i,r_i,s_i',d_i) \sim \mathcal{D}$}\;
        \hlif{#1}{3}{Set $y_i = r_i + \gamma\,(1-d_i)\, Q_{\phi^{-}}\!\big(s_i',\, \mu_{\theta^{-}}(s_i')\big)$}\;
        \hlif{#1}{3}{Update critic by descending $\nabla_\phi \tfrac{1}{N}\sum_i \big(y_i - Q_\phi(s_i,a_i)\big)^2$}\;
        \hlif{#1}{4}{Update actor by ascending $\nabla_\theta \tfrac{1}{N}\sum_i Q_\phi\big(s_i,\, \mu_\theta(s_i)\big)$}\;
        \hlif{#1}{1}{Soft-update targets: $\phi^{-} \leftarrow \tau\phi + (1-\tau)\phi^{-}$; \; $\theta^{-} \leftarrow \tau\theta + (1-\tau)\theta^{-}$}\;
    }
}
\end{algorithm}}

% Caption text reused under each copy of the listing.
\newcommand{\ddpgalgofoot}{{\tiny In code the actor update is a single statement: maximise $\tfrac{1}{N}\sum_i Q_\phi(s_i,\mu_\theta(s_i))$ with $\phi$ held frozen. Autograd's chain rule then recovers the deterministic policy gradient, $\nabla_\theta \mu_\theta \cdot \nabla_a Q_\phi$.}}

\begin{frame}{The DDPG Algorithm}
\vspace{0.1em}
\centering
\ddpgalgo{0}
\end{frame}

\begin{frame}{The DDPG Algorithm --- The Two DQN Tricks}
\begin{columns}[T]
\begin{column}{0.21\textwidth}
\vspace{3.3em}
\footnotesize
\textbf{\textcolor{red}{Replay buffer}}\\[0.4em]
{\scriptsize Collect once, learn from it many times.}
\\[7.5em]
\textbf{\textcolor{red}{``Soft'' target network update}}\\[0.4em]
{\scriptsize Not DQN's hard copy every $C$ steps.}
\end{column}
\begin{column}{0.79\textwidth}
\ddpgalgo{1}
\end{column}
\end{columns}
\end{frame}

% ---------------------------------------------------------------------------
\begin{frame}{Method --- ``Soft'' Target Network Update}
\vspace{0.5em}
Day 2's DQN copies the online weights into the target network \emph{periodically} --- a hard update every $C$ steps. DDPG instead lets the targets \emph{slowly track} the online networks at every step:
\[
\phi^{-} \leftarrow \tau\, \phi + (1-\tau)\, \phi^{-}
\qquad\qquad
\theta^{-} \leftarrow \tau\, \theta + (1-\tau)\, \theta^{-}
\]
\vspace{0.2em}
\begin{itemize}
    \item $\tau \ll 1$ (typically $\tau = 10^{-3}$): the targets change very slowly --- an exponential moving average of the online weights rather than a snapshot of them.
    \pause
    \item \textbf{The cost:} value information propagates backwards through time more slowly. This is the usual stability-versus-speed trade.
\end{itemize}
\end{frame}

% ---------------------------------------------------------------------------
\begin{frame}{The DDPG Algorithm --- Exploration}
\begin{columns}[T]
\begin{column}{0.21\textwidth}
\vspace{5.7em}
\footnotesize
\textbf{\textcolor{red}{Add noise for exploration}}\\[0.4em]
{\scriptsize A deterministic actor explores nothing on its own.}
\end{column}
\begin{column}{0.79\textwidth}
\ddpgalgo{2}
\end{column}
\end{columns}
\end{frame}

% ---------------------------------------------------------------------------
\begin{frame}{Method --- Exploration with a Deterministic Actor}
\vspace{0.4em}
A deterministic actor has \emph{no} exploration of its own: the same state always yields the same action, forever. On Day 3--4 the policy was stochastic, so exploration came for free. Here we must add it by hand:
\[
a_t = \mu_\theta(s_t) + \mathcal{N}_t, \qquad \mathcal{N}_t \sim \mathcal{N}(0, \sigma^2)
\]
\begin{itemize}
    \item Zero-mean \textbf{Gaussian} noise added to the action, with $\sigma$ annealed over training --- the continuous-action analogue of annealing $\varepsilon$ in $\varepsilon$-greedy (Day 1).
    \pause
    \item \textbf{This is exploration only in the behaviour policy.} The stored actions include the noise; the actor being trained does not. Being off-policy is what allows that split.
    \item \textbf{It is also DDPG's weakest point.} The noise is undirected --- it is a random jiggle, not a search. In sparse-reward tasks the agent may never stumble on a reward at all.
\end{itemize}
\end{frame}

% ---------------------------------------------------------------------------
\begin{frame}{The DDPG Algorithm --- Value Network Update}
\begin{columns}[T]
\begin{column}{0.21\textwidth}
\vspace{8.6em}
\footnotesize
\textbf{\textcolor{red}{Value network update}}\\[0.4em]
{\scriptsize Squared TD error --- the same critic loss as Day 2 and Day 4.}
\end{column}
\begin{column}{0.79\textwidth}
\ddpgalgo{3}
\end{column}
\end{columns}
\end{frame}

\begin{frame}{The DDPG Algorithm --- Policy Network Update}
\begin{columns}[T]
\begin{column}{0.21\textwidth}
\vspace{11.5em}
\footnotesize
\textbf{\textcolor{red}{Policy network update}}\\[0.4em]
{\scriptsize Climb the critic. This one line is the deterministic policy gradient.}
\end{column}
\begin{column}{0.79\textwidth}
\ddpgalgo{4}
\end{column}
\end{columns}
\vspace{0.2em}
\ddpgalgofoot
\end{frame}

% ---------------------------------------------------------------------------
% Replaces the two DQN-pseudocode screenshots. Students saw DQN in full on
% Day 2; reprinting the paper's algorithm made this look like a new topic. The
% only thing that actually needed saying is the one-line difference in the
% target, which is what this slide shows.
\begin{frame}{One Line Separates DDPG from DQN}
\vspace{0.6em}
Both algorithms are the same off-policy loop: act, store in $\mathcal{D}$, sample a minibatch, regress $Q$ onto a bootstrapped target. They differ in \emph{how the next action in the target is chosen}.
\vspace{0.8em}
\begin{center}
\renewcommand{\arraystretch}{2.1}
\begin{tabular}{p{0.16\textwidth} p{0.52\textwidth} p{0.22\textwidth}}
\hline
 & \textbf{TD target} & \textbf{Next action from} \\
\hline
\textbf{DQN} (Day 2) & $y = r + \gamma\, \displaystyle\max_{a'} \hat{Q}(s',a';\theta^{-})$ & a \emph{search} over $\mathcal{A}$ \\
\textbf{DDPG} (today) & $y = r + \gamma\, Q_{\phi^{-}}\!\big(s',\, \rbox{\mu_{\theta^{-}}(s')}\big)$ & a \emph{forward pass} \\
\hline
\end{tabular}
\end{center}
\vspace{0.8em}
\begin{itemize}
    \item \textcolor{red}{\textbf{The target actor is the $\arg\max$.}} That single substitution is what carries DQN into continuous action spaces --- and it is why we needed the deterministic policy gradient: something has to \emph{train} $\mu_\theta$ to be that $\arg\max$.
    \item Everything else --- replay buffer, target networks, squared TD error --- is unchanged from Day 2.
\end{itemize}
\end{frame}
